\documentclass[11pt,oneside]{article}

\usepackage[a4paper,margin=25mm,headheight=15pt]{geometry}
\usepackage[T1]{fontenc}
\usepackage[utf8]{inputenc}
\IfFileExists{libertinus.sty}{\usepackage{libertinus}}{\usepackage{lmodern}}
\usepackage{microtype}
\usepackage{amsmath,amssymb,amsthm,mathtools,mathrsfs}
\usepackage{bm}
\usepackage{booktabs,tabularx,array,longtable}
\usepackage{xcolor}
\usepackage{enumitem}
\usepackage{fancyhdr}
\usepackage{titlesec}
\usepackage{listings}
\usepackage{xurl}
\usepackage{hyperref}
\usepackage[nameinlink,noabbrev]{cleveref}

\definecolor{linkblue}{RGB}{25,75,135}
\definecolor{shadegray}{RGB}{246,247,249}
\definecolor{rulegray}{RGB}{195,201,208}
\hypersetup{
  colorlinks=true,
  linkcolor=linkblue,
  citecolor=linkblue,
  urlcolor=linkblue,
  pdftitle={The Fourier Decay of Rvachev's Up-Function: Final Resolution},
  pdfsubject={Dyadic sinc products, lacunary products, transfer operators, Cesaro normalization, and extremal envelopes},
  pdfkeywords={Rvachev up-function, Fabius function, Fourier transform, sinc product, Thue-Morse, transfer operator, Gelfond exponent, Lambert W}
}
\urlstyle{same}

\pagestyle{fancy}
\fancyhf{}
\fancyhead[L]{\small Fourier decay of Rvachev's up-function}
\fancyhead[R]{\small\nouppercase{\leftmark}}
\fancyfoot[C]{\thepage}
% Canonical project-wide mathematical notation.
% Canonical mathematical notation for active FabiusFunction LaTeX documents.
%
% This file is intentionally a plain TeX input rather than a LaTeX package.
% Every standalone document loads it by an explicit path relative to that
% document's directory; included fragments inherit the definitions from their
% root document.  Keep this file independent of document layout and metadata.
%
% Contract:
%   * one command has one mathematical meaning throughout the active corpus;
%   * names encode semantic roles rather than a document-local choice of letter;
%   * visually similar but differently normalized objects have different print;
%   * every command below is defined normatively in the notation catalogue.
%
% The active corpus excludes docs/archive/.  Historical sources may retain
% their original notation only inside an explicitly labelled quotation.

\makeatletter
\@ifundefined{mathbb}{%
  \PackageError{fabius-notation}{The amssymb package must be loaded first}%
    {Load amsmath and amssymb before inputting fabius-notation.tex.}%
}{}
\@ifundefined{operatorname}{%
  \PackageError{fabius-notation}{The amsmath package must be loaded first}%
    {Load amsmath before inputting fabius-notation.tex.}%
}{}
\makeatother

% Version stamp used by the catalogue and by static audits.
\newcommand{\FabiusNotationVersion}{2026-08-31}

% Number systems.  NaturalNumbers is explicitly the nonnegative convention.
\newcommand{\NaturalNumbers}{\mathbb{N}_{0}}
\newcommand{\PositiveIntegers}{\mathbb{N}_{>0}}
\newcommand{\IntegerNumbers}{\mathbb{Z}}
\newcommand{\RationalNumbers}{\mathbb{Q}}
\newcommand{\RealNumbers}{\mathbb{R}}
\newcommand{\ComplexNumbers}{\mathbb{C}}
\newcommand{\ExtendedRealNumbers}{\overline{\mathbb{R}}}

% The three principal Fabius--Rvachev functions.  Subscripts are deliberate:
% a bare F and a calligraphic F have several unrelated uses in the corpus.
\newcommand{\FabiusBounded}{F_{\mathrm{bd}}}
\newcommand{\FabiusGlobal}{F_{\mathrm{gl}}}
\newcommand{\FabiusQuantile}{Q_{F}}
\newcommand{\FabiusClampedQuantile}{Q_{F,\mathrm{cl}}}
\newcommand{\RvachevUp}{\operatorname{up}}

% Constants, elementary operators, and delimiters.
\newcommand{\EulerE}{\mathrm{e}}
\newcommand{\ImaginaryUnit}{\mathrm{i}}
\newcommand{\Differential}{\mathop{}\!\mathrm{d}}
\newcommand{\AbsoluteValue}[1]{\left\lvert #1\right\rvert}
\newcommand{\Norm}[1]{\left\lVert #1\right\rVert}
\newcommand{\Floor}[1]{\left\lfloor #1\right\rfloor}
\newcommand{\Ceiling}[1]{\left\lceil #1\right\rceil}
\newcommand{\FractionalPart}[1]{\left\{#1\right\}}
\newcommand{\PositivePart}[1]{\left(#1\right)_{+}}
\newcommand{\IndicatorOf}[1]{\mathbf{1}_{#1}}
\newcommand{\SetOf}[1]{\left\{#1\right\}}
\newcommand{\InnerProduct}[1]{\left\langle #1\right\rangle}
\newcommand{\CardinalityOf}[1]{\left\lvert #1\right\rvert}
\newcommand{\DefinitionEquals}{\mathrel{\coloneqq}}
\newcommand{\Sign}{\operatorname{sgn}}
\newcommand{\IdentityOperator}{\operatorname{id}}
\newcommand{\SpanOperator}{\operatorname{span}}
\newcommand{\TraceOperator}{\operatorname{tr}}
\newcommand{\RankOperator}{\operatorname{rank}}
\newcommand{\DiagonalOperator}{\operatorname{diag}}
\newcommand{\ResidueOperator}{\operatorname*{Res}}
\newcommand{\LeastCommonMultiple}{\operatorname{lcm}}
\newcommand{\OrderOperator}{\operatorname{ord}}
\newcommand{\PrincipalLogarithm}{\operatorname{Log}}
\newcommand{\FallingFactorial}[2]{\left(#1\right)^{\underline{#2}}}
\newcommand{\RisingFactorial}[2]{\left(#1\right)^{\overline{#2}}}

% Support, topology, convolution, and metric notation.
\newcommand{\SupportOperator}{\operatorname{supp}}
\newcommand{\SupportOf}[1]{\SupportOperator\!\left(#1\right)}
\newcommand{\TopologicalSupportOf}[1]{\operatorname{tsupp}\!\left(#1\right)}
\newcommand{\Convolution}{\mathbin{*}}
\newcommand{\MetricDistance}{\operatorname{dist}}
\newcommand{\TotalVariationDistance}{d_{\mathrm{TV}}}

% Probability.  Equality and convergence in law are not metric distance.
\newcommand{\Probability}{\mathbb{P}}
\newcommand{\Expectation}{\mathbb{E}}
\newcommand{\Variance}{\operatorname{Var}}
\newcommand{\Covariance}{\operatorname{Cov}}
\newcommand{\LawOperator}{\operatorname{Law}}
\newcommand{\LawOf}[1]{\LawOperator\!\left(#1\right)}
\newcommand{\EqualInLaw}{\mathrel{\overset{\mathrm{law}}{=}}}
\newcommand{\ConvergesInLaw}{\mathrel{\xrightarrow{\mathrm{law}}}}

% Fourier and integral transforms.  The printed labels expose normalization.
% FourierTwoPi uses exp(-2*pi*i*x*xi); FourierAngular uses exp(-i*x*omega).
\newcommand{\FourierTwoPi}[1]{\widehat{#1}_{\,2\pi}}
\newcommand{\FourierAngular}[1]{\widehat{#1}_{\,\mathrm{ang}}}
\newcommand{\LaplaceTransformSymbol}{\mathcal{L}}
\newcommand{\MellinTransformSymbol}{\mathcal{M}}
\newcommand{\LaplaceTransformOf}[1]{\LaplaceTransformSymbol\!\left[#1\right]}
\newcommand{\MellinTransformOf}[1]{\MellinTransformSymbol\!\left[#1\right]}
\newcommand{\MomentGeneratingFunction}[1]{M_{#1}}
\newcommand{\CharacteristicFunction}[1]{\varphi_{#1}}

% The two standard sinc functions must never share one symbol.
\newcommand{\SincRad}{\operatorname{sinc}_{\mathrm{rad}}}
\newcommand{\SincPi}{\operatorname{sinc}_{\pi}}

% Binary arithmetic and Thue--Morse notation.
\newcommand{\BinaryDigitSum}{\operatorname{wt}_{2}}
\newcommand{\ThueMorseSignSymbol}{\varepsilon}
\newcommand{\ThueMorseSign}[1]{\varepsilon_{#1}}
\newcommand{\PadicValuation}[1]{\nu_{#1}}
\newcommand{\TwoAdicValuation}{\nu_{2}}

% q-series notation.  Every base is an explicit argument.
\newcommand{\QPochhammer}[3]{\left(#1;#2\right)_{#3}}
\newcommand{\GaussianBinomial}[3]{%
  \genfrac{[}{]}{0pt}{}{#1}{#2}_{#3}%
}
\newcommand{\QInteger}[2]{\left[#1\right]_{#2}}

% Named special functions.
\newcommand{\Polylogarithm}{\operatorname{Li}}
\newcommand{\LambertW}{W}
\newcommand{\GammaFunction}{\Gamma}
\newcommand{\RiemannZeta}{\zeta}

% Asymptotics.  BigO and LittleO are operator tokens for existing formula
% grammar; prefer the At forms so that the limiting regime is printed.
\newcommand{\BigO}{\mathcal{O}}
\newcommand{\LittleO}{o}
\newcommand{\BigOAt}[2]{\mathcal{O}_{#2}\!\left(#1\right)}
\newcommand{\LittleOAt}[2]{o_{#2}\!\left(#1\right)}

\renewcommand{\headrulewidth}{0.4pt}

\titleformat{\section}{\Large\bfseries}{\thesection}{0.7em}{}
\titleformat{\subsection}{\large\bfseries}{\thesubsection}{0.7em}{}
\titleformat{\subsubsection}{\normalsize\bfseries}{\thesubsubsection}{0.7em}{}
\setcounter{secnumdepth}{3}
\setcounter{tocdepth}{2}
\setlist[itemize]{leftmargin=2em,itemsep=0.25em,topsep=0.4em}
\setlist[enumerate]{leftmargin=2.3em,itemsep=0.25em,topsep=0.4em}
\setlength{\emergencystretch}{3em}
\allowdisplaybreaks[2]

\newtheorem{theorem}{Theorem}[section]
\newtheorem{proposition}[theorem]{Proposition}
\newtheorem{lemma}[theorem]{Lemma}
\newtheorem{corollary}[theorem]{Corollary}
\newtheorem{conjecture}[theorem]{Conjecture}
\theoremstyle{definition}
\newtheorem{definition}[theorem]{Definition}
\newtheorem{example}[theorem]{Example}
\theoremstyle{remark}
\newtheorem{remark}[theorem]{Remark}
\newtheorem{warning}[theorem]{Warning}

\newcommand{\spr}{\operatorname{spr}}
\newcommand{\T}{\mathsf T}
\newcommand{\dist}{\mathop{\longrightarrow}\limits^{\mathrm d}}
\newcommand{\LambertW}{W}

\newenvironment{boxedremark}{%
  \par\smallskip\noindent\begin{center}\begin{minipage}{0.95\textwidth}%
  \hrule\vspace{0.65em}\small
}{%
  \vspace{0.65em}\hrule\end{minipage}\end{center}\smallskip
}

\lstdefinestyle{code}{
  basicstyle=\ttfamily\footnotesize,
  backgroundcolor=\color{shadegray},
  frame=single,
  rulecolor=\color{rulegray},
  columns=fullflexible,
  keepspaces=true,
  showstringspaces=false,
  breaklines=true,
  xleftmargin=0.6em,
  xrightmargin=0.6em,
  aboveskip=0.8em,
  belowskip=0.8em
}

\title{\bfseries The Fourier Decay of Rvachev's Up-Function\\[0.35em]
\large A final resolution of the pointwise, metric, averaged, and extremal regimes}
\author{Prepared for Vladimir Reshetnikov's \texttt{ProveIt} development}
\date{26 August 2026}

\begin{document}
\maketitle

\begin{abstract}
Let
\[
 f(x)=\widehat{\RvachevUp}(x)
 =\prod_{j=0}^{\infty}\SincRad\!\left(\frac{\pi x}{2^j}\right),
 \qquad \SincRad z=\frac{\sin z}{z},
\]
where the Fourier convention is
$\widehat h(\xi)=\int_{\RealNumbers}h(t)\EulerE^{-2\pi\ImaginaryUnit t\xi}\Differential t$.
The motivating question asks for the real-axis decay of $f$, and in particular
for a positive analytic, eventually decreasing gauge $g$ for which the
ordinary Ces\`aro mean of $|f|/g$ either converges to one or remains bounded
between two positive constants.  Four first-wave reports and a later
comparative audit produced several apparently competing answers.

The competition disappears once the word ``amplitude'' is separated into its
inequivalent meanings.  The exact identity
\[
 |f(2^k(1+t))|
 =\frac{2^{-k(k+1)/2}}{\pi^{k+1}(1+t)^{k+1}}
 H(1+t)\prod_{j=0}^{k}|\sin(\pi2^jt)|,
 \qquad 0\le t\le1,
\]
reduces the problem to a lacunary product for the doubling map.  Every natural
size has the common log-normal core
\[
 \exp\!\left(-\frac{(\log x)^2}{2\log2}\right),
\]
but its power of $x$ depends on whether one takes a geometric mean, an
arithmetic mean, an RMS mean, or a supremum.  The respective exponents are
\[
\begin{aligned}
 \kappa_0&=3.151496129472319\ldots,
 &\kappa_1&=2.748070014871334\ldots,\\
 \kappa_2&=2.651496129472318\ldots,
 &\kappa_\infty&=2.359014879111741\ldots .
\end{aligned}
\]
The arithmetic exponent is governed by the Perron root
$\rho_1=0.661322602060565\ldots$ of an explicit transfer operator.  With
\[
 E_\kappa(x)=x^{-\kappa}
 \exp\!\left(-\frac{(\log x)^2}{2\log2}\right)
\]
and
$A_1=0.0912661241315\ldots$, the elementary gauge
$g=A_1E_{\kappa_1}$ is positive, real analytic, eventually decreasing, and
satisfies the requested two-sided Ces\`aro bound.  The normalized mean tends
to one at dyadic endpoints.  For unrestricted $t=2^ky$, however, it converges
to a nonconstant dyadic-mantissa profile determined by a singular conformal
eigenmeasure.  The same obstruction defeats every gauge whose within-shell
shape stabilizes, including the usual elementary, regularly varying, and
continuous log-periodic corrections.  This is the precise negative result;
the materials do not prove nonexistence among arbitrary scale-dependent
analytic gauges.  A logarithmic Ces\`aro average, with measure $\Differential x/x$, does
converge for all $t$ after the explicit rescaling
$A_1^{\log}=0.0938606357568\ldots$.

For almost every fixed dyadic ray the centered logarithm satisfies a CLT and a
two-sided law of the iterated logarithm with asymptotic variance $\pi^2/4$ per
dyadic step.  This corrects a factor $2$ in the variance calculation in the
Aistleitner--Hofer--Larcher source: the sharp coefficient in front of
$\sqrt{n\log\log n}$ is $\pi/\sqrt2$, not $\pi$.

Finally, if
$M_k=\max_{2^k\le x\le2^{k+1}}|f(x)|$, put
$a=\log2$, $\beta=\log(\sqrt3/2)$,
$c=(a/\pi)\sqrt{3/2}$, and $n=k+1$.  Then
\[
 \log M_k=
 -\frac a2k(k+1)-n\log\pi+n\beta
 -\frac{\LambertW(cn)^2+2\LambertW(cn)}{2a}
 +\BigO((\log\log k)^2).
\]
Thus the true annular envelope contains a second log-normal suppression in
$\log\log x$, not a fixed power of $\log x$.  Individual interzero peaks are
still more irregular: immediately after $m2^k$, with $m$ fixed and odd, they
obey an explicit much smaller two-adic asymptotic.  Together these results give
a single consistent resolution of all supplied materials.
\end{abstract}

\tableofcontents
\newpage

\section{The question, the conventions, and the final verdict}
\label{sec:intro}

Rvachev's up-function is an even, nonnegative $C^\infty$ function supported on
$[-1,1]$ and normalized to have integral one.  Under the Fourier convention in
the abstract, its transform is the infinite sinc product above; see
Arias de Reyna~\cite{arias1982}, the Fabius-function summary~\cite{wikipedia},
and the formal and expository development in the \texttt{ProveIt}
repository~\cite{proveit}.

The original question~\cite{mse-source} asks for a positive analytic,
eventually decreasing $g$ such that, for a sufficiently large $t_0$,
\begin{equation}
 \frac1{t-t_0}\int_{t_0}^{t}\frac{|f(x)|}{g(x)}\Differential x
 \longrightarrow1,
 \qquad t\to\infty,
 \label{eq:Q2}
\end{equation}
or at least such that this expression stays between two positive constants.
The latter condition will be denoted by \emph{Q3}; the full limit will be
denoted by \emph{Q2}.

A single pointwise equivalent cannot exist: $f$ vanishes at every nonzero
integer.  More importantly, even after removing those zeros, four different
operations on the same lacunary numerator give four different exponential
rates.  Put
\begin{equation}
 a:=\log2,
 \qquad
 E_\kappa(x):=x^{-\kappa}\exp\!\left(-\frac{(\log x)^2}{2a}\right).
 \label{eq:E-kappa}
\end{equation}
The complete first-order spectrum is as follows.

\begin{center}
\small
\begin{tabularx}{\textwidth}{@{}>{\raggedright\arraybackslash}p{0.25\textwidth}
 >{\centering\arraybackslash}p{0.17\textwidth}
 >{\centering\arraybackslash}p{0.23\textwidth}>{\raggedright\arraybackslash}X@{}}
\toprule
Notion of size & numerator rate $\Lambda$ & exponent
$\displaystyle\kappa=\frac{\log\pi+a/2-\log\Lambda}{a}$ & conclusion\\
\midrule
Geometric mean / a.e. fixed ray
& $1/2$
& $\kappa_0=3.1514961295\ldots$
& Typical logarithmic law, with CLT and LIL fluctuations.\\[0.6em]
Arithmetic mean ($L^1$)
& $\rho_1=0.6613226021\ldots$
& $\kappa_1=2.7480700149\ldots$
& Correct exponent for the Ces\`aro question; Q3 holds.\\[0.6em]
Root mean square ($L^2$)
& $1/\sqrt2$
& $\kappa_2=\log_2(2\pi)=2.6514961295\ldots$
& Exact identity $\int_0^1P_n^2=2^{-n}$.\\[0.6em]
Uniform fixed-ray envelope
& $\sqrt3/2$
& $\kappa_\infty=2.3590148791\ldots$
& Sharp global upper envelope, attained on the $1/3\leftrightarrow2/3$ orbit.\\
\bottomrule
\end{tabularx}
\end{center}

The exponent $\log_2(2\pi)$ proposed in the early draft is therefore not a
random near miss.  It is exactly the RMS exponent.  It is not the pointwise,
arithmetic-mean, or supremum exponent, and no fixed power of $\log x$ can
repair a mismatch already proportional to $\log x$ in the logarithm.

\begin{boxedremark}
\textbf{Final answer to the original normalization question.}
Let
\[
 g_1(x)=A_1E_{\kappa_1}(x),
 \qquad
 A_1=0.0912661241315\ldots .
\]
Then Q3 holds for all sufficiently large $t$.  The Q2 limit equals one along
$t=2^k$, but for $t=2^ky$, $1\le y\le2$, it converges to a nonconstant
continuous profile $\mathscr P(y)$.  Hence this natural elementary gauge does
not satisfy unrestricted Q2.  No gauge with a stabilized continuous
within-shell shape can satisfy Q2.  This includes all standard log-normal
power gauges, slowly varying corrections, and continuous log-periodic
modulations.  The unrestricted class of arbitrary scale-dependent analytic
monotone gauges is not eliminated by the supplied arguments and should not be
claimed to be eliminated.
\end{boxedremark}

\subsection{Division between elementary proof, imported analysis, and numerics}

The exact product identities, the zero and lobe structure, the $L^2$ identity,
the Gelfond bound, the Lambert-$W$ optimization, the singularity argument, and
the two-adic peak asymptotic are proved below directly.  The metric results use
standard ergodic and martingale limit theorems.  The finite-$q$ shell limits use
the standard Ruelle--Perron--Frobenius spectral theorem for a contracting
two-branch weighted transfer operator.  The existence of the $L^1$ exponential
rate also follows from Fouvry--Mauduit~\cite{fouvrymauduit}; the rigorous
published enclosure $0.66130<\rho_1<0.66135$ is due to
Aistleitner--Hofer--Larcher~\cite{ahl2017}.  Decimal digits beyond a certified
enclosure are numerical, not substitutes for proof.

\section{Entire product, zero multiplicities, signs, and lobes}
\label{sec:product}

\subsection{Convergence, scaling, and the canonical zero product}

The infinite product
\begin{equation}
 f(z)=\prod_{j=0}^{\infty}\SincRad\!\left(\frac{\pi z}{2^j}\right)
 \label{eq:f-def}
\end{equation}
converges locally uniformly on $\ComplexNumbers$.  Indeed,
$\log\SincRad w=-w^2/6+\BigO(w^4)$ near zero, and
$\sum_j|z|^2/4^j$ is locally uniformly bounded.  Hence $f$ is even and entire,
$f(0)=1$, and its factors may be regrouped locally uniformly.

\begin{proposition}[Three exact product descriptions]
\label{prop:products}
For every $z\in\ComplexNumbers$,
\begin{align}
 f(2z)&=\SincRad(2\pi z)f(z),
 \label{eq:dilation}\\
 f(z)&=\prod_{m=1}^{\infty}
 \left(1-\frac{z^2}{m^2}\right)^{1+\TwoAdicValuation(m)},
 \label{eq:weierstrass}\\
 f(z)&=\SincRad(\pi z)H(z),
 \qquad
 H(z):=\prod_{j=1}^{\infty}\SincRad\!\left(\frac{\pi z}{2^j}\right).
 \label{eq:H-def}
\end{align}
Thus every nonzero integer $m$ is a zero of multiplicity $1+\TwoAdicValuation(|m|)$,
and there are no other zeros.
\end{proposition}

\begin{proof}
Shifting the index in \eqref{eq:f-def} gives \eqref{eq:dilation}; separating
the first factor gives \eqref{eq:H-def}.  Euler's product
\[
 \SincRad(\pi z)=\prod_{r=1}^{\infty}\left(1-\frac{z^2}{r^2}\right)
\]
applied at $z/2^j$ shows that a fixed integer $m$ appears once for each
$j\in\{0,\ldots,\TwoAdicValuation(m)\}$.  This gives \eqref{eq:weierstrass} and the stated
zero multiplicities.
\end{proof}

\subsection{Exactly one extremum in every interzero interval}

The Weierstrass product exposes more than the zeros.  Away from the integers,
termwise logarithmic differentiation gives
\begin{align}
 (\log|f|)'(x)
 &=-2x\sum_{m=1}^{\infty}
   \frac{1+\TwoAdicValuation(m)}{m^2-x^2},
 \label{eq:logder}\\
 (\log|f|)''(x)
 &=-2\sum_{m=1}^{\infty}(1+\TwoAdicValuation(m))
   \frac{m^2+x^2}{(m^2-x^2)^2}<0.
 \label{eq:logconcave}
\end{align}
The differentiated series converge locally uniformly off the zero set.

\begin{theorem}[Lobe structure and signed Thue--Morse pattern]
\label{thm:lobes}
The function $f$ decreases strictly from $1$ to $0$ on $(0,1)$.  For every
integer $m\ge1$, the interval $(m,m+1)$ contains exactly one critical point
$c_m$; it is the unique maximum of $|f|$ on that interval.  Moreover,
\begin{equation}
 \operatorname{sgn}f(x)=(-1)^{\BinaryDigitSum(m)},
 \qquad m<x<m+1,
 \label{eq:TM-sign}
\end{equation}
where $\BinaryDigitSum(m)$ is the sum of the binary digits of $m$.
\end{theorem}

\begin{proof}
Strict negativity in \eqref{eq:logconcave} makes $(\log|f|)'$ strictly
decreasing on every interzero interval.  On $(m,m+1)$ its pole at the left
endpoint has limit $+\infty$, while its pole at the right endpoint has limit
$-\infty$; hence it has exactly one zero.  On $(0,1)$ it starts at zero and is
strictly decreasing, so $f$ is strictly decreasing there.

Crossing the zero at $r$ changes the sign $1+\TwoAdicValuation(r)$ times.  The sign on
$(m,m+1)$ is therefore
\[
 (-1)^{\sum_{r=1}^{m}(1+\TwoAdicValuation(r))}.
\]
Legendre's identity $\TwoAdicValuation(m!)=m-\BinaryDigitSum(m)$ gives
$\sum_{r=1}^{m}(1+\TwoAdicValuation(r))=2m-\BinaryDigitSum(m)$, whose parity is $\BinaryDigitSum(m)$.
\end{proof}

The theorem legitimizes speaking of ``the peak'' in each unit interval, but it
does not make the peak sequence regular.  In \cref{sec:two-adic} we will see
that the peaks immediately following large powers of two are extraordinarily
deep troughs.

\section{The exact dyadic reduction}
\label{sec:factorization}

Put
\begin{equation}
 s(t):=|\sin(\pi t)|,
 \qquad
 \T t:=2t\pmod1,
 \qquad
 P_n(t):=\prod_{j=0}^{n-1}s(\T^jt).
 \label{eq:Pn}
\end{equation}
The product $P_n$ is, up to $2^{-n}$, the modulus of the length-$2^n$
Thue--Morse trigonometric polynomial:
\begin{equation}
 2^nP_n(t)=
 \left|\prod_{j=0}^{n-1}(1-\EulerE^{2\pi\ImaginaryUnit2^jt})\right|
 =\left|\sum_{r=0}^{2^n-1}(-1)^{\BinaryDigitSum(r)}\EulerE^{2\pi\ImaginaryUnit rt}\right|.
 \label{eq:TM-polynomial}
\end{equation}

\begin{theorem}[Exact shell factorization]
\label{thm:factorization}
Let $k\ge0$, $n=k+1$, $0\le t\le1$, and $y=1+t$.  Then
\begin{equation}
 \boxed{
 |f(2^ky)|
 =2^{-k(k+1)/2}(\pi y)^{-n}H(y)P_n(t).}
 \label{eq:exact-factorization}
\end{equation}
Here $H$ is given by \eqref{eq:H-def}; it is positive on $[1,2)$ and vanishes
at $2$.
\end{theorem}

\begin{proof}
Iterating \eqref{eq:dilation} gives
\[
 f(2^ky)=f(y)\prod_{j=1}^{k}\SincRad(\pi2^jy).
\]
Since $2^j$ is an integer,
$|\sin(\pi2^jy)|=|\sin(\pi2^jt)|$.  The $k$ new denominators multiply to
$\pi^ky^k2^{k(k+1)/2}$, while
$|f(y)|=H(y)|\sin(\pi t)|/(\pi y)$.  Collecting the $n=k+1$ sine factors
proves the formula.
\end{proof}

\subsection{One algebraic normalization for all notions of size}

For a candidate numerator growth rate $0<\Lambda<1$, define
\begin{equation}
 \kappa(\Lambda):=
 \frac{\log\pi+a/2-\log\Lambda}{a}
 \label{eq:kappa-Lambda}
\end{equation}
and
\begin{equation}
 W_\Lambda(y):=
 H(y)\frac{\Lambda}{\pi}
 y^{\kappa(\Lambda)-1}
 \exp\!\left(\frac{(\log y)^2}{2a}\right),
 \qquad 1\le y\le2.
 \label{eq:W-Lambda}
\end{equation}
The function $W_\Lambda$ is continuous, nonnegative, bounded, and positive on
$[1,2)$.

\begin{corollary}[Master normalization identity]
\label{cor:master}
With the notation of \cref{thm:factorization},
\begin{equation}
 \boxed{
 \frac{|f(2^ky)|}{E_{\kappa(\Lambda)}(2^ky)}
 =\Lambda^{-n}W_\Lambda(y)P_n(t).}
 \label{eq:master}
\end{equation}
\end{corollary}

\begin{proof}
Write $b=\log y$ and $\log(2^ky)=ka+b$.  The logarithm of the deterministic
factor in \eqref{eq:exact-factorization}, divided by $E_\kappa$, is
\[
 k(\kappa a-\log\pi-a/2)
 +\frac{b^2}{2a}+(\kappa-1)b-\log\pi.
\]
For \eqref{eq:kappa-Lambda}, the coefficient of $k$ is $-\log\Lambda$.
Moving one further factor $\Lambda^{-1}$ into the bounded mantissa term gives
exactly \eqref{eq:W-Lambda} and \eqref{eq:master}.
\end{proof}

The triangular denominator $2^{-k(k+1)/2}$ is responsible for the universal
quadratic term in $\log x$.  Every disagreement in the lower-order term is now
localized in the operation applied to $P_n$.

\section{Typical dyadic rays and the corrected metric fluctuation law}
\label{sec:metric}

\subsection{Birkhoff averages and the typical exponent}

Set
\begin{equation}
 \psi(t):=\log|2\sin(\pi t)|=a+\log s(t).
 \label{eq:psi}
\end{equation}
This function belongs to $L^p(0,1)$ for every finite $p$, and
\begin{equation}
 \int_0^1\psi(t)\Differential t=0.
 \label{eq:psi-mean}
\end{equation}
For $S_n(t):=\sum_{j=0}^{n-1}\psi(\T^jt)$, the master identity with
$\Lambda=1/2$ becomes the exact equality
\begin{equation}
 \log|f(2^k(1+t))|
 +\frac{(\log(2^k(1+t)))^2}{2a}
 +\kappa_0\log(2^k(1+t))
 =\log W_{1/2}(1+t)+S_n(t),
 \label{eq:exact-centered}
\end{equation}
where
\begin{equation}
 \kappa_0=\kappa(1/2)
 =\frac32+\frac{\log\pi}{a}
 =3.151496129472319\ldots .
 \label{eq:kappa0}
\end{equation}

\begin{theorem}[Typical ray law]
\label{thm:typical}
For Lebesgue-almost every $t\in(0,1)$,
\begin{equation}
 \log|f(2^k(1+t))|
 =-\frac{(\log(2^k(1+t)))^2}{2a}
  -\kappa_0\log(2^k(1+t))+o(k).
 \label{eq:typical}
\end{equation}
\end{theorem}

\begin{proof}
The doubling map is ergodic for Lebesgue measure.  Birkhoff's theorem applied
to the integrable centered function \eqref{eq:psi} gives $S_n(t)=o(n)$ almost
everywhere.  Substitute this into \eqref{eq:exact-centered}.
\end{proof}

More generally, whenever the orbit average
$\alpha(t)=\lim n^{-1}\sum_{j<n}\log s(\T^jt)$ exists and is finite, the same
algebra gives the ray exponent
\begin{equation}
 \kappa(\alpha)=\frac{\log\pi+a/2-\alpha}{a}.
 \label{eq:kappa-alpha}
\end{equation}
Thus arithmetic properties of the binary expansion remain visible on
exceptional rays.

\subsection{An exact reverse-martingale decomposition}

Let $P$ be the Perron operator of the doubling map,
\begin{equation}
 (Ph)(x)=\frac12\left(h(x/2)+h((x+1)/2)\right).
 \label{eq:Perron-unweighted}
\end{equation}
The double-angle identity gives
\begin{equation}
 P\psi=\frac12\psi.
 \label{eq:Ppsi}
\end{equation}
Define
\begin{equation}
 d:=2\psi-\psi\circ\T.
 \label{eq:d-def}
\end{equation}
Then
\begin{equation}
 d(x)=\log|\tan(\pi x)|,
 \qquad Pd=0,
 \label{eq:d-tan}
\end{equation}
and therefore $d\circ\T^j$ is a stationary reverse martingale-difference
sequence.  Telescoping gives
\begin{equation}
 S_n(t)=\sum_{j=0}^{n-1}d(\T^jt)-\psi(t)+\psi(\T^nt).
 \label{eq:martingale-coboundary}
\end{equation}

\begin{lemma}[Exact asymptotic variance]
\label{lem:variance}
One has
\begin{equation}
 \int_0^1d(x)^2\Differential x=\frac{\pi^2}{4}.
 \label{eq:d-variance}
\end{equation}
Equivalently,
\begin{equation}
 \int_0^1\psi(x)\psi(\T^rx)\Differential x
 =\frac{\pi^2}{12\,2^r},
 \qquad r\ge0,
 \label{eq:covariance}
\end{equation}
and hence
\begin{equation}
 \Variance(S_n)
 =\frac{\pi^2}{12}\left(3n-4+2^{2-n}\right)
 =\frac{\pi^2}{4}n+\BigO(1).
 \label{eq:finite-variance}
\end{equation}
\end{lemma}

\begin{proof}
By symmetry and the standard beta-integral derivatives,
\[
 \int_0^1\log^2|\tan(\pi x)|\Differential x
 =\frac2\pi\int_0^{\pi/2}\log^2(\tan u)\Differential u
 =\frac{\pi^2}{4}.
\]
For the Fourier proof, use
\begin{equation}
 \psi(x)=-\sum_{m=1}^{\infty}\frac{\cos(2\pi mx)}m
 \quad\text{in }L^2(0,1).
 \label{eq:psi-Fourier}
\end{equation}
Orthogonality contributes the factor $1/2$ and gives
\[
 \int\psi(x)\psi(2^rx)\Differential x
 =\frac12\sum_{m=1}^{\infty}\frac1{(2^rm)m}
 =\frac{\pi^2}{12\,2^r}.
\]
Summing the covariance series yields \eqref{eq:finite-variance}.
\end{proof}

\begin{theorem}[CLT and two-sided LIL]
\label{thm:CLT-LIL}
If $t$ is Lebesgue distributed on $(0,1)$, then
\begin{equation}
 \frac{S_n(t)}{(\pi/2)\sqrt n}\dist\mathcal N(0,1).
 \label{eq:CLT}
\end{equation}
For Lebesgue-almost every fixed $t$,
\begin{align}
 \limsup_{n\to\infty}
 \frac{S_n(t)}{(\pi/2)\sqrt{2n\log\log n}}&=1,
 \label{eq:LIL-plus}\\
 \liminf_{n\to\infty}
 \frac{S_n(t)}{(\pi/2)\sqrt{2n\log\log n}}&=-1.
 \label{eq:LIL-minus}
\end{align}
\end{theorem}

\begin{proof}
The variable $d$ has finite moments of every order.  Since $Pd=0$, the
martingale CLT and martingale LIL apply to the first term on the right of
\eqref{eq:martingale-coboundary}, with variance \eqref{eq:d-variance}; see
Hall--Heyde~\cite{hallheyde} and Stout~\cite{stout}.  The logarithmic
singularities of $\psi$ have exponentially decreasing tails, so
$\psi(\T^nt)=o(\sqrt{n\log\log n})$ almost surely by Borel--Cantelli.  The
fixed term $\psi(t)$ is harmless.
\end{proof}

In the original frequency variable, the almost-sure upper and lower
log-amplitude excursions have size
\begin{equation}
 \pm\frac{\pi}{\sqrt{2a}}
 \sqrt{\log x\,\log\log\log x}
 \quad\text{along a typical dyadic ray}.
 \label{eq:LIL-x}
\end{equation}
They are larger than every fixed multiple of $\log\log x$, so no factor
$(\log x)^p$ can make the typical normalized amplitude converge.

\subsection{Correction to the published lacunary-product constant}
\label{subsec:AHL-correction}

Aistleitner--Hofer--Larcher~\cite{ahl2017} state, in their Theorem~2, a
coefficient $\pi$ in front of $\sqrt{n\log\log n}$ for the upper excursions of
$\prod_{j<n}|2\sin(\pi2^jt)|$.  Their displayed variance calculation writes
\[
 \sum_{m\ge1}\left(\frac1{m^2}
 +2\sum_{r\ge1}\frac1{2^rm^2}\right)=\frac{\pi^2}{2}.
\]
The cosine orthogonality factor $1/2$ is missing.  Formula
\eqref{eq:covariance} shows that the correct Green--Kubo variance is
\begin{equation}
 \frac12\sum_{m\ge1}\frac1{m^2}
 +2\sum_{r\ge1}\frac12\sum_{m\ge1}\frac1{2^rm^2}
 =\frac{\pi^2}{4}.
 \label{eq:AHL-corrected-variance}
\end{equation}
Consequently the correct coefficient relative to
$\sqrt{n\log\log n}$ is
\begin{equation}
 \boxed{\frac{\pi}{\sqrt2}},
 \label{eq:AHL-corrected-constant}
\end{equation}
not $\pi$.  The exact martingale decomposition above supplies an independent
proof and also yields the missing lower LIL.  The third first-wave report
inherited the published coefficient; the first report and the comparative
audit correctly detected the variance issue.

\section{The complete finite-\texorpdfstring{$L^q$}{Lq} shell spectrum}
\label{sec:Lq}

\subsection{The weighted doubling operator}

For $q>0$, define
\begin{equation}
 (\LaplaceTransformSymbol_qh)(x):=\frac12\left[
 \sin^q\!\left(\frac{\pi x}{2}\right)h(x/2)
 +\cos^q\!\left(\frac{\pi x}{2}\right)h((x+1)/2)
 \right].
 \label{eq:Lq-operator}
\end{equation}
This is the Perron operator of the doubling map with weight $s^q$.  A change
of variables on the two inverse branches gives the exact duality formula
\begin{equation}
 \int_0^1 h(t)P_n(t)^q\Differential t
 =\int_0^1\LaplaceTransformSymbol_q^nh(x)\Differential x.
 \label{eq:Lq-duality}
\end{equation}

The inverse branches contract by $1/2$, and the weights are H\"older of every
order not exceeding $\min(1,q)$.  The endpoint zeros do not destroy
primitivity: at $x=0$ and $x=1$, the surviving branch has positive weight.
The standard Ruelle--Perron--Frobenius theorem therefore gives the following;
see Baladi~\cite{baladi} or Ruelle~\cite{ruelle}.

\begin{proposition}[Spectral data]
\label{prop:RPF}
For each $q>0$, the operator $\LaplaceTransformSymbol_q$ on a suitable H\"older space has a
simple leading eigenvalue $\rho_q>0$, a strictly positive eigenfunction $h_q$,
and a positive eigenmeasure $\nu_q$.  They may be normalized by
\begin{equation}
 \LaplaceTransformSymbol_qh_q=\rho_qh_q,
 \qquad
 \LaplaceTransformSymbol_q^*\nu_q=\rho_q\nu_q,
 \qquad
 \nu_q(h_q)=1.
 \label{eq:RPF-normalization}
\end{equation}
There is $0<\vartheta_q<1$ such that
\begin{equation}
 \rho_q^{-n}\LaplaceTransformSymbol_q^nh
 =h_q\nu_q(h)+\BigO(\vartheta_q^n\Norm h_{C^\alpha})
 \label{eq:RPF-gap}
\end{equation}
for compatible H\"older $h$.
\end{proposition}

Set
\begin{equation}
 \Lambda_q:=\rho_q^{1/q},
 \qquad
 \kappa_q:=\kappa(\Lambda_q)
 =\frac{\log\pi+a/2-\log\Lambda_q}{a}.
 \label{eq:kappa-q}
\end{equation}
For a function $F$ on the $k$th shell, write
\begin{equation}
 \mathcal M_{q,k}(F):=
 \left(2^{-k}\int_{2^k}^{2^{k+1}}|F(x)|^q\Differential x\right)^{1/q}.
 \label{eq:shell-Lq}
\end{equation}

\begin{theorem}[Sharp finite-$q$ shell normalization]
\label{thm:Lq-spectrum}
For every $q\in(0,\infty)$ there is $A_q\in(0,\infty)$ such that
\begin{equation}
 \mathcal M_{q,k}\!\left(\frac{f}{E_{\kappa_q}}\right)
 \longrightarrow A_q,
 \qquad k\to\infty.
 \label{eq:Lq-limit}
\end{equation}
More precisely,
\begin{equation}
 A_q^q=
 \left(\int_0^1h_q(x)\Differential x\right)
 \nu_q\!\left(W_{\Lambda_q}(1+\cdot)^q\right).
 \label{eq:Aq}
\end{equation}
The quantity $\Lambda_q$ is nondecreasing in $q$, and $\kappa_q$ is
nonincreasing.
\end{theorem}

\begin{proof}
Use $x=2^k(1+t)$ in \eqref{eq:shell-Lq} and raise
\eqref{eq:master} to the $q$th power.  Since
$\Lambda_q^{-nq}=\rho_q^{-n}$,
\[
 \mathcal M_{q,k}\!\left(\frac{f}{E_{\kappa_q}}\right)^q
 =\rho_q^{-n}\int_0^1W_{\Lambda_q}(1+t)^qP_n(t)^q\Differential t.
\]
Apply \eqref{eq:Lq-duality} and \eqref{eq:RPF-gap}.  This gives
\eqref{eq:Lq-limit} and \eqref{eq:Aq}.  For each fixed $n$, the $L^q$ norm of
$P_n$ on the probability space $[0,1]$ is nondecreasing in $q$.  Take $n$th
roots and then $n\to\infty$ to obtain monotonicity of $\Lambda_q$; formula
\eqref{eq:kappa-q} reverses the order for $\kappa_q$.
\end{proof}

A factor $(\log x)^p$ is not part of any sharp finite-$q$ normalization.  At
the correct $\kappa_q$, inserting such a factor multiplies the normalized
shell mean by $(k\log2)^{-p}(1+o(1))$; it tends to zero for $p>0$ and infinity
for $p<0$.

\subsection{The four distinguished values}

\paragraph{Geometric mean, $q=0$.}
Because Lebesgue measure is invariant under $\T$ and
$\int_0^1\log s=-a$, the limiting numerator rate is $1/2$.  In fact, the
normalized shell geometric mean at $\kappa_0$ is exactly independent of $k$:
\begin{equation}
 \exp\!\left(2^{-k}\int_{2^k}^{2^{k+1}}
 \log\frac{|f(x)|}{E_{\kappa_0}(x)}\Differential x\right)
 =\exp\!\left(\int_0^1\log W_{1/2}(1+t)\Differential t\right).
 \label{eq:geometric-exact}
\end{equation}
This is the shell form of \cref{thm:typical}.

\paragraph{Arithmetic mean, $q=1$.}
Write
\begin{equation}
 \rho:=\rho_1.
 \label{eq:rho}
\end{equation}
Fouvry and Mauduit proved
\begin{equation}
 \int_0^1P_n(t)\Differential t=K\rho^n(1+o(1)),
 \qquad K>0,
 \label{eq:FM}
\end{equation}
and Aistleitner--Hofer--Larcher established
\begin{equation}
 0.66130<\rho<0.66135.
 \label{eq:rho-certified}
\end{equation}
High-precision transfer-operator computation gives
\begin{equation}
 \rho=0.661322602060565\ldots .
 \label{eq:rho-numeric}
\end{equation}
Thus
\begin{equation}
 \boxed{
 \kappa_1=\frac12+\frac{\log(\pi/\rho)}a
 =2.748070014871334\ldots .}
 \label{eq:kappa1}
\end{equation}
This is the exponent relevant to the ordinary absolute mean.

\paragraph{Root mean square, $q=2$.}
Since
\begin{equation}
 \LaplaceTransformSymbol_2\mathbf1
 =\frac12\left(\sin^2\frac{\pi x}{2}
 +\cos^2\frac{\pi x}{2}\right)=\frac12\mathbf1,
 \label{eq:L2-one}
\end{equation}
we have $\rho_2=1/2$ and
\begin{equation}
 \int_0^1P_n(t)^2\Differential t=2^{-n}
 \label{eq:L2-exact}
\end{equation}
for every $n$.  Consequently
\begin{equation}
 \boxed{\kappa_2=\log_2(2\pi)=2.651496129472318\ldots .}
 \label{eq:kappa2}
\end{equation}
This identifies precisely what the original draft computed.

\paragraph{Uniform norm, $q=\infty$.}
The limit is
$\Lambda_\infty=\sqrt3/2$ and gives
\begin{equation}
 \boxed{
 \kappa_\infty
 =\frac32+\frac{\log(\pi/\sqrt3)}a
 =2.359014879111741\ldots .}
 \label{eq:kappa-infty}
\end{equation}
A direct sharp proof is given in \cref{sec:extremal}.

The strict ordering
\begin{equation}
 \kappa_\infty<\kappa_2<\kappa_1<\kappa_0
 \label{eq:kappa-order}
\end{equation}
is the cleanest summary of why the four reports appeared to disagree.

\section{The ordinary Ces\`aro question}
\label{sec:Cesaro}

\subsection{The sharp arithmetic gauge and Q3}

Let
\begin{equation}
 W_1(y):=W_\rho(y)
 =H(y)\frac{\rho}{\pi}
 y^{\kappa_1-1}
 \exp\!\left(\frac{(\log y)^2}{2a}\right).
 \label{eq:W1}
\end{equation}
Define
\begin{equation}
 A_1:=
 \left(\int_0^1h_1(x)\Differential x\right)
 \nu_1\!\left(W_1(1+\cdot)\right)
 =0.0912661241315\ldots .
 \label{eq:A1}
\end{equation}
By \cref{thm:Lq-spectrum},
\begin{equation}
 2^{-k}\int_{2^k}^{2^{k+1}}
 \frac{|f(x)|}{E_{\kappa_1}(x)}\Differential x
 \longrightarrow A_1.
 \label{eq:L1-shell-limit}
\end{equation}

\begin{theorem}[Answer to Q3 and dyadic Q2]
\label{thm:Q3}
Put
\begin{equation}
 g_1(x):=A_1E_{\kappa_1}(x).
 \label{eq:g1}
\end{equation}
For every fixed $t_0>0$, there are $0<c<C<\infty$ such that
\begin{equation}
 c\le
 \frac1{t-t_0}\int_{t_0}^{t}\frac{|f(x)|}{g_1(x)}\Differential x
 \le C
 \label{eq:Q3-result}
\end{equation}
for all sufficiently large $t$.  Moreover,
\begin{equation}
 \frac1{2^K-t_0}\int_{t_0}^{2^K}
 \frac{|f(x)|}{g_1(x)}\Differential x\longrightarrow1.
 \label{eq:dyadic-Q2}
\end{equation}
\end{theorem}

\begin{proof}
Let $J_k$ be the integral of $|f|/E_{\kappa_1}$ over
$[2^k,2^{k+1}]$.  Equation \eqref{eq:L1-shell-limit} says
$J_k=A_1 2^k(1+o(1))$.  Summing the geometric sequence of completed shells
proves \eqref{eq:dyadic-Q2}.  If $2^K\le t<2^{K+1}$, the preceding complete
shell already contributes a fixed positive multiple of $2^K$, while the sum
of all completed shells and the whole current shell is at most a fixed
multiple of $2^K$.  Since $t\asymp2^K$, this proves \eqref{eq:Q3-result}.
\end{proof}

\subsection{Analyticity and eventual derivative signs}

The requested regularity is not an issue for \eqref{eq:g1}.

\begin{proposition}[Eventual alternating derivatives]
\label{prop:derivatives}
The function $g_1$ is positive and real analytic on $(0,\infty)$.  For each
fixed integer $m\ge0$, there exists $X_m$ such that
\begin{equation}
 (-1)^m g_1^{(m)}(x)>0,
 \qquad x\ge X_m.
 \label{eq:alternating-derivatives}
\end{equation}
In particular it is eventually strictly decreasing, and every fixed
derivative is eventually monotone.
\end{proposition}

\begin{proof}
Put $u(x)=\log x/a+\kappa_1$.  Then
$g_1'(x)=-x^{-1}u(x)g_1(x)$.  Induction using
$\Differential/\Differential x=x^{-1}\Differential/\Differential\log x$ gives
\begin{equation}
 \frac{g_1^{(m)}(x)}{g_1(x)}
 =x^{-m}\left((-1)^m u(x)^m+\BigO_m(u(x)^{m-1})\right).
 \label{eq:derivative-asymptotic}
\end{equation}
The leading term controls the sign for sufficiently large $x$.
\end{proof}

The threshold $X_m$ depends on $m$; this is not a claim of complete
monotonicity on one common half-line.

\subsection{The nonconstant dyadic-mantissa profile}

The shell limit is stronger than mere boundedness but weaker than unrestricted
Q2.  Normalize $h_1,\nu_1$ as in \eqref{eq:RPF-normalization} and define
\begin{equation}
 \mathscr F(z):=
 \frac{\nu_1(W_1(1+\cdot)\mathbf1_{[0,z]})}
      {\nu_1(W_1(1+\cdot))},
 \qquad 0\le z\le1.
 \label{eq:F-profile}
\end{equation}

\begin{lemma}[The conformal measure is atomless and singular]
\label{lem:singular}
The eigenmeasure $\nu_1$ has no atoms and is singular with respect to Lebesgue
measure.  The same is true after multiplication by the continuous weight
$W_1(1+z)$.
\end{lemma}

\begin{proof}
For an atom at $z$, the eigenmeasure equation gives
\begin{equation}
 \rho\,\nu_1(\{z\})
 =\frac12s(z)\nu_1(\{\T z\}).
 \label{eq:atom-relation}
\end{equation}
Endpoint atoms vanish because $s=0$.  If a remaining atom existed, its forward
masses would grow by at least the factor $2\rho>1$ until an orbit repeated; a
periodic orbit gives an immediate contradiction and a nonperiodic orbit gives
infinitely many atoms of unbounded mass.  Hence $\nu_1$ is atomless.

Suppose its absolutely continuous component has a nonzero density $r\in L^1$.
The Lebesgue decomposition is preserved by the two nonsingular inverse
branches, so that component separately satisfies
\begin{equation}
 \rho r(x)=s(x)r(\T x)
 \quad\text{for a.e. }x.
 \label{eq:conformal-density}
\end{equation}
Iteration gives
\begin{equation}
 r(\T^nx)=\frac{\rho^nr(x)}{P_n(x)}.
 \label{eq:density-iteration}
\end{equation}
For Lebesgue-a.e. $x$, Birkhoff's theorem gives
$P_n(x)=2^{-n+o(n)}$, so the right side grows like
$(2\rho)^{n+o(n)}r(x)$.  On the other hand, for every $\varepsilon>0$,
Markov's inequality and Borel--Cantelli imply
$r(\T^nx)\le\EulerE^{\varepsilon n}$ eventually for a.e. $x$, because $\T$
preserves Lebesgue measure.  Taking
$0<\varepsilon<\log(2\rho)$ is a contradiction on the set $r>0$.  Thus
$\nu_1$ is singular.  Multiplying by a continuous positive weight away from a
single atomless endpoint preserves singularity and absence of atoms.
\end{proof}

Because the weighted measure is atomless, $\mathscr F$ is continuous.

\begin{theorem}[Full dyadic phase limit]
\label{thm:phase}
For every $y\in[1,2]$ and every fixed $t_0>0$,
\begin{equation}
 \lim_{K\to\infty}
 \frac1{2^Ky-t_0}\int_{t_0}^{2^Ky}
 \frac{|f(x)|}{A_1E_{\kappa_1}(x)}\Differential x
 =\mathscr P(y):=\frac{1+\mathscr F(y-1)}{y}.
 \label{eq:phase-limit}
\end{equation}
The function $\mathscr P$ is continuous,
$\mathscr P(1)=\mathscr P(2)=1$, and it is not identically one.
Consequently $g_1$ does not satisfy unrestricted Q2.
\end{theorem}

\begin{proof}
From \eqref{eq:master} and $x=2^K(1+t)$,
\begin{equation}
 2^{-K}\int_{2^K}^{2^K(1+z)}
 \frac{|f(x)|}{E_{\kappa_1}(x)}\Differential x
 =\rho^{-n}\int_0^zW_1(1+t)P_n(t)\Differential t.
 \label{eq:partial-shell}
\end{equation}
The spectral decomposition \eqref{eq:RPF-gap}, followed by approximation of
$\mathbf1_{[0,z]}$ from above and below by H\"older functions, shows that the
right side tends to $A_1\mathscr F(z)$.  All preceding complete shells
contribute $A_12^K(1+o(1))$.  Division by $A_12^Ky$ proves
\eqref{eq:phase-limit}.

If $\mathscr P\equiv1$, then
$1+\mathscr F(y-1)=y$, or $\mathscr F(z)=z$.  This would make the normalized
weighted eigenmeasure Lebesgue measure, contradicting \cref{lem:singular}.
\end{proof}

\subsection{Exactly what kind of Q2 gauge is ruled out}

The singular profile yields a robust impossibility theorem, but its hypotheses
matter.

\begin{theorem}[No stabilized within-shell gauge]
\label{thm:no-stabilized-gauge}
Suppose $g$ is positive for large $x$ and that there exist positive constants
$c_K$ and a positive continuous function $u$ on $[0,1]$ such that
\begin{equation}
 \frac{g(2^K(1+z))}
      {c_KE_{\kappa_1}(2^K(1+z))}
 \longrightarrow u(z)
 \label{eq:stabilized-gauge}
\end{equation}
uniformly for $0\le z\le1$.  Then $g$ cannot satisfy Q2.
\end{theorem}

\begin{proof}
Under \eqref{eq:stabilized-gauge}, the normalized partial-shell measures are,
up to $1/c_K+o(1)$,
\[
 \rho^{-n}\frac{W_1(1+z)}{u(z)}P_n(z)\Differential z.
\]
The spectral theorem makes these converge, after total-mass normalization, to
the probability measure proportional to
$W_1(1+z)u(z)^{-1}\Differential\nu_1(z)$.  Q2 would force every partial shell to have
asymptotic mass equal to its Euclidean length, hence would force that limiting
measure to be Lebesgue.  It is singular by \cref{lem:singular}, a
contradiction.
\end{proof}

Slowly varying corrections satisfy \eqref{eq:stabilized-gauge} with $u=1$.
A continuous log-periodic multiplier gives a continuous positive $u$.  The
same applies to the standard elementary and Hardy-field log-normal gauges once
the correct quadratic and linear terms have been fixed.  What the theorem
does \emph{not} address is a deliberately scale-dependent analytic gauge whose
shape develops ever finer structure from shell to shell.  The comparative
audit's phrase ``unattainable in principle'' is therefore too broad if read as
a theorem about every conceivable analytic monotone function.

\subsection{A logarithmic Ces\`aro mean removes the phase}

The phase survives ordinary averaging because the newest dyadic shell has the
same order as all preceding shells combined.  With the measure $\Differential x/x$,
every shell has equal logarithmic length and the final incomplete shell is
negligible.

Define
\begin{equation}
 B_1:=\lim_{k\to\infty}
 \int_{2^k}^{2^{k+1}}
 \frac{|f(x)|}{E_{\kappa_1}(x)}\frac{\Differential x}{x}
 =0.0650592350404\ldots
 \label{eq:B1log}
\end{equation}
and
\begin{equation}
 A_1^{\log}:=\frac{B_1}{a}
 =0.0938606357568\ldots .
 \label{eq:A1log}
\end{equation}
The existence of the limit follows from the same spectral argument, with
$W_1(1+t)$ replaced by $W_1(1+t)/(1+t)$.

\begin{corollary}[All-phase logarithmic Ces\`aro limit]
\label{cor:log-Cesaro}
For every fixed $t_0>0$,
\begin{equation}
 \frac1{\log(t/t_0)}
 \int_{t_0}^{t}
 \frac{|f(x)|}{A_1^{\log}E_{\kappa_1}(x)}\frac{\Differential x}{x}
 \longrightarrow1.
 \label{eq:log-Cesaro}
\end{equation}
\end{corollary}

\begin{proof}
Every complete dyadic shell contributes $B_1+o(1)$ to the unscaled numerator,
whereas its logarithmic length is $a$.  The number of complete shells is
asymptotic to $\log t/a$.  The one incomplete shell contributes only
$\BigO(1)$ and is negligible after division by $\log t$.
\end{proof}

\section{The sharp fixed-ray and uniform envelope}
\label{sec:extremal}

The $L^\infty$ numerator rate can be proved by one elementary inequality.  Put
\begin{equation}
 \lambda:=\frac{\sqrt3}{2},
 \qquad
 \beta:=\log\lambda.
 \label{eq:lambda-beta}
\end{equation}

\begin{lemma}[Sharp two-step sine inequality]
\label{lem:two-step}
For every real $x$,
\begin{equation}
 \frac23\log s(x)+\frac13\log s(\T x)
 \le\log\lambda,
 \label{eq:two-step-log}
\end{equation}
with the usual value $-\infty$ at a zero.  Equality holds exactly at
$x\equiv1/3$ or $2/3\pmod1$.
\end{lemma}

\begin{proof}
Cubing the exponential form reduces the assertion to
$s(x)^2s(\T x)\le\lambda^3$.  Write
$u=\sin^2(\pi x)$.  The square of the left side is
$4u^3(1-u)$, whose maximum on $[0,1]$ is $27/64$, attained only at
$u=3/4$.  Taking square roots gives
$3\sqrt3/8=\lambda^3$ and identifies the equality orbit.
\end{proof}

\begin{proposition}[Exact exponential rate of the finite supremum]
\label{prop:Pn-sup}
For every $n\ge1$,
\begin{equation}
 \lambda^n\le\Norm{P_n}_\infty\le\lambda^{n-1}.
 \label{eq:Pn-sup}
\end{equation}
Hence
\begin{equation}
 \lim_{n\to\infty}\Norm{P_n}_\infty^{1/n}=\lambda.
 \label{eq:Gelfond-rate}
\end{equation}
\end{proposition}

\begin{proof}
At $x=1/3$, the orbit alternates between $1/3$ and $2/3$, so every factor is
$\lambda$.  For the upper bound, sum \eqref{eq:two-step-log} over
$x,\T x,\ldots,\T^{n-2}x$.  If $\phi=\log s$, the result is
\[
 2\phi(x)+3\sum_{j=1}^{n-2}\phi(\T^jx)
 +\phi(\T^{n-1}x)\le3(n-1)\log\lambda.
\]
Adding the nonpositive boundary combination
$\phi(x)+2\phi(\T^{n-1}x)$ shows
$3\sum_{j<n}\phi(\T^jx)\le3(n-1)\log\lambda$.
Exponentiate.
\end{proof}

This is the classical Gelfond exponent for the Thue--Morse polynomial; the
more general ergodic-optimization setting is discussed by Fan, Schmeling, and
Shen~\cite{fanschmeling}.

\begin{theorem}[Sharp global envelope]
\label{thm:uniform-envelope}
With $\kappa_\infty$ from \eqref{eq:kappa-infty}, there is $C<\infty$ such
that
\begin{equation}
 |f(x)|\le C E_{\kappa_\infty}(x),
 \qquad x\ge1.
 \label{eq:uniform-envelope}
\end{equation}
The exponent is sharp.  More precisely, for every $k\ge0$,
\begin{equation}
 |f(2^k\tfrac43)|
 =W_\lambda(\tfrac43)E_{\kappa_\infty}(2^k\tfrac43).
 \label{eq:exact-max-ray}
\end{equation}
Consequently
\begin{equation}
 \limsup_{x\to\infty}
 \frac{\log|f(x)|+(\log x)^2/(2a)}{\log x}
 =-\kappa_\infty.
 \label{eq:global-limsup}
\end{equation}
\end{theorem}

\begin{proof}
Use \eqref{eq:master} with $\Lambda=\lambda$ and
\eqref{eq:Pn-sup}:
\[
 \frac{|f(2^k(1+t))|}{E_{\kappa_\infty}(2^k(1+t))}
 \le\lambda^{-1}\max_{1\le y\le2}W_\lambda(y).
\]
At $t=1/3$, $P_n(t)=\lambda^n$ exactly, which proves
\eqref{eq:exact-max-ray}.  Taking logarithms gives
\eqref{eq:global-limsup}.
\end{proof}

The exact ray \eqref{eq:exact-max-ray} is sharp relative to its \emph{actual}
frequency.  It is not the largest value in the shell whose left endpoint is
$2^k$: the fixed mantissa $4/3$ pays a power-sized displacement penalty.  The
true shell maximizer starts much closer to the left endpoint and spends a
slowly growing transition time reaching the maximizing two-cycle.  That is the
next problem.

\section{The sharp dyadic-annulus maximum and Lambert \texorpdfstring{$W$}{W}}
\label{sec:annular}

Define
\begin{equation}
 M_k:=\max_{2^k\le x\le2^{k+1}}|f(x)|,
 \qquad n:=k+1.
 \label{eq:Mk}
\end{equation}
From \eqref{eq:exact-factorization},
\begin{align}
 M_k&=\pi^{-n}2^{-k(k+1)/2}A_n,
 \label{eq:Mk-An}\\
 A_n&:=\max_{0\le t\le1}
 H(1+t)(1+t)^{-n}P_n(t).
 \label{eq:An}
\end{align}
All nontrivial lower-order behavior lies in $A_n$.

\subsection{The boundary-layer coordinates}

For $0<t\le1/2$, write
\begin{equation}
 t=2^{-r}s,
 \qquad r\in\IntegerNumbers_{\ge0},
 \qquad \frac14\le s\le\frac12.
 \label{eq:transition-coordinates}
\end{equation}
When $r<n$,
\begin{equation}
 P_n(2^{-r}s)
 =\left(\prod_{h=1}^{r}\sin\frac{\pi s}{2^h}\right)P_{n-r}(s).
 \label{eq:P-transition}
\end{equation}
The small-angle block has the uniform factorization
\begin{equation}
 \prod_{h=1}^{r}\sin\frac{\pi s}{2^h}
 =(\pi s)^r2^{-r(r+1)/2}H_r(s),
 \qquad
 H_r(s):=\prod_{h=1}^{r}\SincRad\frac{\pi s}{2^h}.
 \label{eq:small-block}
\end{equation}
On the compact interval $[1/4,1/2]$, $H_r$ is bounded above and below by
positive constants uniformly in $r$.

The upper tail estimate follows from \cref{prop:Pn-sup}.  A matching lower
estimate, with enough freedom to optimize over $s$, comes from dense preimages
of the maximizing cycle.

\begin{lemma}[Near-maximizing preimages]
\label{lem:preimages}
There is an absolute $C>0$ such that, for every $s_0\in[1/4,1/2]$ and every
integer $q\ge1$, one can choose $s_q\in[1/4,1/2]$ with
\begin{equation}
 |s_q-s_0|\le2^{-q},
 \qquad
 \T^qs_q\in\{1/3,2/3\},
 \label{eq:preimage-close}
\end{equation}
and, for every $m\ge q$,
\begin{equation}
 P_m(s_q)\ge\exp(m\beta-Cq^2).
 \label{eq:preimage-lower}
\end{equation}
\end{lemma}

\begin{proof}
The $q$fold preimages of $1/3$ are
$(j+1/3)/2^q$ and have spacing $2^{-q}$; using the preimages of $2/3$ near an
endpoint gives \eqref{eq:preimage-close}.  Before time $q$, every orbit point
is a nonintegral rational with reduced denominator dividing
$3\cdot2^{q-j}$.  Since
$|\sin\pi x|\ge2\operatorname{dist}(x,\IntegerNumbers)$, the first $q$ factors have
logarithm bounded below by $-Cq^2$.  Thereafter every factor is
$\sqrt3/2=\EulerE^\beta$.
\end{proof}

Introduce the finite-dimensional objective
\begin{equation}
 \DecayOptimizationObjective{n}(r,s):=
 -\frac a2r^2
 +r\left(\log(\pi s)-\frac a2-\beta\right)
 -n\log(1+s2^{-r}).
 \label{eq:Fn}
\end{equation}

\begin{lemma}[Reduction to the boundary objective]
\label{lem:An-reduction}
As $n\to\infty$,
\begin{equation}
 \log A_n
 =n\beta+
 \sup_{\substack{r\in\IntegerNumbers_{\ge0}\\1/4\le s\le1/2}}
 \DecayOptimizationObjective{n}(r,s)
 +\BigO((\log\log n)^2).
 \label{eq:An-reduction}
\end{equation}
\end{lemma}

\begin{proof}
For the upper bound, the region $t\ge1/2$ loses a fixed exponential factor
through $(1+t)^{-n}$ and is below the eventual benchmark.  In the region
$0<t<1/2$, use \eqref{eq:P-transition}, \eqref{eq:small-block}, and
$P_{n-r}(s)\le\lambda^{n-r-1}$.  The factors $H(1+t)$ and $H_r(s)$ contribute
only $\BigO(1)$.  Values $r\ge n$ have all factors in the small-angle regime
and are of order $\exp(-an^2/2+\BigO(n))$, hence are negligible.  This proves
the claimed upper bound.

For the lower bound, take $(r,s_0)$ within one of the supremum.  Comparison
with $r\asymp\log n$ shows that every such near-maximizer satisfies
\begin{equation}
 r=\BigO(\log n),
 \qquad
 ns_02^{-r}=\BigO((\log n)^2).
 \label{eq:localization}
\end{equation}
Choose
$q=\lceil4\log_2\log(n+\EulerE^\EulerE)\rceil$ and replace $s_0$ by the point $s_q$
from \cref{lem:preimages}.  On the localized set, the $s$ derivative of
$\DecayOptimizationObjective{n}$ is $\BigO((\log n)^2)$, so this replacement changes the
objective by $\BigO(1)$.  The tail loses only
$Cq^2=\BigO((\log\log n)^2)$ in \eqref{eq:preimage-lower}.  All remaining
factors are bounded below uniformly.  This proves the matching lower bound.
\end{proof}

\subsection{The exact Lambert optimization}

Put
\begin{equation}
 C_0:=\log\pi-\frac a2-\beta
 =\log\!\left(\pi\sqrt{\frac23}\right),
 \qquad
 c:=a\EulerE^{-C_0}=\frac a\pi\sqrt{\frac32}.
 \label{eq:c-Lambert}
\end{equation}
Numerically, $c=0.270222319733365\ldots$.

\begin{lemma}[Lambert-$W$ saddle]
\label{lem:Lambert}
Let $w_n=\LambertW(cn)$, on the principal branch.  Then
\begin{equation}
 \sup_{r,s}\DecayOptimizationObjective{n}(r,s)
 =-\frac{w_n^2+2w_n}{2a}+\BigO(1),
 \label{eq:Lambert-result}
\end{equation}
where the supremum is over the set in \eqref{eq:An-reduction}.
\end{lemma}

\begin{proof}
At a maximizer, \eqref{eq:localization} implies that replacing
$-n\log(1+s2^{-r})$ by $-ns2^{-r}$ changes the value by $o(1)$.  Let
$A=C_0+\log s$ and $z=ar-A$.  A direct completion of the square gives
\begin{equation}
 -\frac a2r^2+r(C_0+\log s)-ns2^{-r}
 =-\frac{z^2}{2a}-n\EulerE^{-C_0}\EulerE^{-z}+\frac{A^2}{2a}.
 \label{eq:z-objective}
\end{equation}
As $s$ ranges over $[1/4,1/2]$, $A$ ranges over an interval of length exactly
$a$.  Successive integer values of $r$ therefore make the corresponding
$z$ intervals meet end to end and tile the real line.  The last term in
\eqref{eq:z-objective} is uniformly bounded, so the supremum differs by
$\BigO(1)$ from the real maximum of
\[
 -\frac{z^2}{2a}-n\EulerE^{-C_0}\EulerE^{-z}.
\]
Its critical equation is $z\EulerE^z=an\EulerE^{-C_0}=cn$, hence $z=w_n$.  At this
point $n\EulerE^{-C_0}\EulerE^{-w_n}=w_n/a$, which gives
\eqref{eq:Lambert-result}.
\end{proof}

The tiling is important: the continuous mantissa $s$, whose permitted
multiplicative width is exactly two, absorbs the integer rounding in $r$.
Without it one would lose an avoidable $\BigO(\log n)$ term.

\begin{theorem}[Sharp dyadic-annulus envelope]
\label{thm:annular-main}
Let $n=k+1$, $w_k=\LambertW(cn)$, and
\begin{equation}
 \gamma:=\frac a2+\log\pi-\beta
 =\log\!\left(\frac{2^{3/2}\pi}{\sqrt3}\right).
 \label{eq:gamma}
\end{equation}
Then
\begin{equation}
 \boxed{
 \log M_k=
 -\frac a2k(k+1)-n\log\pi+n\beta
 -\frac{w_k^2+2w_k}{2a}
 +\BigO((\log\log k)^2).}
 \label{eq:annular-centered}
\end{equation}
Equivalently,
\begin{equation}
 \log M_k=
 -\frac a2k^2-\gamma k
 -\frac{\LambertW(c(k+1))^2+2\LambertW(c(k+1))}{2a}
 +\BigO((\log\log k)^2).
 \label{eq:annular-compact}
\end{equation}
\end{theorem}

\begin{proof}
Combine \eqref{eq:Mk-An}, \cref{lem:An-reduction}, and
\cref{lem:Lambert}.  The difference between the two displayed forms is a
constant absorbed by the remainder.
\end{proof}

Expanding
$\LambertW(cn)=\log n-\log\log n+\log c+\BigO((\log\log n)/\log n)$ gives
\begin{align}
 \log M_k={}&-\frac a2k^2-\gamma k
 -\frac{(\log k)^2}{2a}
 +\frac{\log k\log\log k}{a}
 \notag\\
 &-\frac{1+\log c}{a}\log k
 +\BigO((\log\log k)^2).
 \label{eq:annular-expanded}
\end{align}

In the physical scale $R=2^k$, put $L=\log R$ and
$d=\pi^{-1}\sqrt{3/2}$.  Since $c(k+1)=d(L+a)$,
\begin{equation}
 \begin{split}
 \log\max_{R\le x\le2R}|f(x)|
 ={}&-\frac{L^2}{2a}-\kappa_\infty L\\
 &-\frac{\LambertW(d(L+a))^2+2\LambertW(d(L+a))}{2a}
 +\BigO((\log\log L)^2).
 \end{split}
 \label{eq:annular-R}
\end{equation}
The first terms of the extra boundary loss are
\begin{equation}
 -\frac{(\log\log R)^2}{2a}
 +\frac{\log\log R\,\log\log\log R}{a}
 +\BigO(\log\log R).
 \label{eq:second-lognormal}
\end{equation}
This is smaller than every fixed negative power of $\log R$, yet larger than
every fixed negative power of $R$.  It is the decisive correction missed by a
normalizer containing only $(\log R)^{-p}$.

The proof also localizes a near-maximizer.  At the smooth saddle,
\begin{equation}
 t\asymp\frac{\LambertW(cn)}{an},
 \qquad
 r=\frac{\LambertW(cn)}a+\BigO(1)
 =\log_2n-\log_2\log n+\BigO(1).
 \label{eq:max-location}
\end{equation}
Thus the largest shell peak lies at
$x=2^k(1+\Theta(\log k/k))$, not at a fixed mantissa.

\section{Exceptional two-adic troughs among the local peaks}
\label{sec:two-adic}

Let
\begin{equation}
 E_m:=\max_{m\le x\le m+1}|f(x)|=|f(c_m)|,
 \label{eq:Em}
\end{equation}
where $c_m$ is the unique critical point from \cref{thm:lobes}.  The sequence
$E_m$ has no single smooth asymptotic, even within one dyadic shell.

Fix an odd positive integer $m$, put $N=m2^k$, and write $x=N+t$ with
$0\le t\le1$.  Define
\begin{equation}
 f_k(t):=\prod_{j=0}^{k}\SincRad\!\left(\frac{\pi t}{2^j}\right).
 \label{eq:fk}
\end{equation}
For $0\le j\le k$, $N/2^j$ is an integer and
\begin{equation}
 \left|\SincRad\!\left(\frac{\pi(N+t)}{2^j}\right)\right|
 =\frac{t}{N+t}\SincRad\!\left(\frac{\pi t}{2^j}\right).
\end{equation}
The remaining tail gives the exact identity
\begin{equation}
 |f(N+t)|
 =\left(\frac{t}{N+t}\right)^{k+1}
 f_k(t)\left|H\!\left(m+\frac{t}{2^k}\right)\right|.
 \label{eq:near-m2k}
\end{equation}
Because $m$ is odd, $H(m)\ne0$, and uniformly on $[0,1]$,
\begin{equation}
 f_k(t)=f(t)(1+\BigO(4^{-k})),
 \qquad
 H\!\left(m+\frac{t}{2^k}\right)=H(m)(1+\BigO(2^{-k})).
 \label{eq:tail-uniform}
\end{equation}

\begin{lemma}[A linear endpoint zero against a large power]
\label{lem:boundary-max}
Let $F$ be positive on $[0,1)$ and
$F(t)=h(1-t)+\BigO((1-t)^2)$ as $t\uparrow1$, with $h>0$.  If
$J_p=\max_{0\le t\le1}t^pF(t)$, then
\begin{equation}
 J_p\sim\frac{h}{\EulerE p},
 \qquad
 p(1-t_p)\longrightarrow1
 \label{eq:boundary-max}
\end{equation}
for every maximizing point $t_p$.
\end{lemma}

\begin{proof}
Set $s=p(1-t)$.  Uniformly for bounded $s$,
$p t^pF(t)\to hs\EulerE^{-s}$, whose unique maximum is $h/\EulerE$ at $s=1$.
The bound $t^pF(t)\le C(1-t)\EulerE^{-p(1-t)}$ near $1$, together with exponential
smallness away from $1$, prevents a maximizing $s$ from escaping to zero or
infinity.
\end{proof}

Since $(\SincRad(\pi t))'|_{t=1}=-1$, formula \eqref{eq:H-def} gives
\begin{equation}
 f(t)=H(1)(1-t)+\BigO((1-t)^2),
 \text{ as }t\uparrow1,
 \qquad
 H(1)=0.553771275888810\ldots .
 \label{eq:f-near-one}
\end{equation}

\begin{theorem}[Peaks following an odd multiple of a power of two]
\label{thm:two-adic-peaks}
Let $m$ be fixed and odd, and write $c_{m2^k}=m2^k+t_k$.  Then
\begin{equation}
 (k+1)(1-t_k)\longrightarrow1
 \label{eq:peak-location-two-adic}
\end{equation}
and
\begin{equation}
 \boxed{
 E_{m2^k}\sim
 \frac{H(1)|H(m)|}{\EulerE(k+1)}(m2^k)^{-(k+1)}.}
 \label{eq:two-adic-peak}
\end{equation}
In particular,
\begin{equation}
 E_{2^k}\sim
 \frac{H(1)^2}{\EulerE(k+1)}2^{-k(k+1)}.
 \label{eq:power-two-peak}
\end{equation}
\end{theorem}

\begin{proof}
Put $p=k+1$.  Equations \eqref{eq:near-m2k} and \eqref{eq:tail-uniform} give,
uniformly in $t$,
\[
 |f(m2^k+t)|
 =(m2^k)^{-p}|H(m)|t^pf(t)(1+o(1)).
\]
Apply \cref{lem:boundary-max} and \eqref{eq:f-near-one}.  The uniform
multiplicative error preserves both the value and location asymptotics.
\end{proof}

Comparing \eqref{eq:power-two-peak} with \eqref{eq:annular-compact} gives
\begin{equation}
 \log\frac{E_{2^k}}{M_k}=-\frac a2k^2+\BigO(k).
 \label{eq:peak-spread}
\end{equation}
Thus peaks whose locations differ by less than a factor two can differ by an
additional log-normal factor.  This decisively rules out a smooth asymptotic
curve for all interzero maxima.

\section{Comparative audit of the supplied materials}
\label{sec:audit}

The reports are best viewed as studying different functionals of the same
finite product $P_n$, not as proposing mutually exclusive pointwise formulas.
The following table records the substantive contribution and the correction
needed for each source.

\begin{center}
\small
\begin{tabularx}{\textwidth}{@{}>{\raggedright\arraybackslash}p{0.18\textwidth}
 >{\raggedright\arraybackslash}p{0.39\textwidth}>{\raggedright\arraybackslash}X@{}}
\toprule
Material & Correct contribution & Required correction or scope restriction\\
\midrule
Original MSE source~\cite{mse-source}
& Correctly asks an averaged question and already notices integer zeros,
Thue--Morse signs, and a log-normal-looking scale.
& It does not specify which mean of the lacunary numerator determines the
power correction; the supplied answer does not solve Q2 or Q3.\\[0.7em]

First report~\cite{wave1a}
& Exact product and lobe structure; typical exponent; exact martingale
coboundary; CLT and the corrected LIL variance; sharp annular expansion through
$\BigO(\log k)$.
& It does not solve the $L^1$/Ces\`aro problem, and its annular error is weaker
than the later Lambert-$W$ result.\\[0.7em]

Second report~\cite{wave1b}
& Strongest treatment of the dyadic-annulus maximum: exact Lambert-$W$
formula with $\BigO((\log\log k)^2)$ error; dense-preimage construction; exact
two-adic peak troughs.
& It explicitly recognizes that the original Ces\`aro question is a separate
$L^1$ problem and therefore does not answer it.\\[0.7em]

Third report~\cite{wave1c}
& Exact dyadic factorization; typical, $L^1$, $L^2$, and $L^\infty$ exponents;
Q3; dyadic endpoint normalization; singular dyadic profile.
& It imports the coefficient $\pi$ from Aistleitner--Hofer--Larcher.  The
correct coefficient is $\pi/\sqrt2$ relative to
$\sqrt{n\log\log n}$.\\[0.7em]

Fourth report~\cite{wave1d}
& Most complete transfer-operator spectrum; exact $L^2$ identity; explicit
Collatz--Wielandt certificate for $\rho_1$; Q3 and phase profile; sharp uniform
envelope; unique lobes and exceptional two-adic peaks.
& Its uniform envelope is a fixed-ray/global-limsup statement, not the finer
maximum relative to the left endpoint of each dyadic shell.  The Lambert
boundary layer from the second report is still needed.\\[0.7em]

AHL lacunary-product source~\cite{ahl2017}
& Supplies the relevant lacunary-product framework and a rigorous enclosure
for the $L^1$ growth constant.
& Its displayed $f_1$ variance omits the cosine Parseval factor $1/2$, doubling
the variance and multiplying the LIL coefficient by $\sqrt2$.\\[0.7em]

Comparative audit~\cite{comparative}
& Correctly unifies the $L^q$ spectrum, recognizes the LIL variance error,
derives Q3 and the singular phase obstruction, and adds logarithmic averaging.
& ``Unattainable in principle'' is justified only for gauges with stabilized
within-shell shape.  Its quoted numerical phase span comes from a coarse
mantissa grid and misses narrower features of the singular-continuous profile.\\
\bottomrule
\end{tabularx}
\end{center}

\subsection{The two apparent contradictions that are not contradictions}

\paragraph{Fixed maximizing ray versus shell maximum.}
Along $x=2^k(4/3)$, the ratio $|f(x)|/E_{\kappa_\infty}(x)$ is exactly
constant.  In the shell $[2^k,2^{k+1}]$, however, comparing with the scale at
the left endpoint $2^k$ charges the fixed ray the factor
$E_{\kappa_\infty}(4\cdot2^k/3)/E_{\kappa_\infty}(2^k)$, which is a negative
power of $2^k$.  A boundary-layer point much closer to $2^k$ is larger even
though it has not followed the maximizing orbit from time zero.  Hence the
exact fixed-ray theorem and the extra Lambert suppression in the shell maximum
are simultaneously true.

\paragraph{The RMS exponent versus the arithmetic exponent.}
The identity $\int P_n^2=2^{-n}$ makes $1/\sqrt2$ the exact RMS numerator rate.
The arithmetic numerator rate is the different Perron root
$\rho=0.6613226\ldots$.  Since $\rho<1/\sqrt2$, the arithmetic-normalized
amplitude decays with the larger power $\kappa_1>\kappa_2$.  Treating an
$L^2$ computation as a pointwise or $L^1$ computation is exactly the source of
the early exponent mismatch.

\subsection{What is now final, and what remains genuinely open}

The following statements are settled by the arguments assembled here.
\begin{enumerate}
 \item The exact dyadic reduction and all four first-order exponents.
 \item The typical CLT and corrected two-sided LIL.
 \item Q3 for an explicit elementary analytic gauge and Q2 at dyadic endpoints.
 \item Failure of full Q2 for that gauge and for every stabilized within-shell
 gauge.
 \item The sharp fixed-ray/global envelope and the sharper dyadic-annulus
 Lambert-$W$ asymptotic.
 \item The exceptional two-adic asymptotic for local peaks.
\end{enumerate}
The only normalization issue not settled by the supplied methods is the bare
existence of a deliberately non-self-similar, scale-dependent, positive
analytic monotone gauge satisfying ordinary Q2.  Such a gauge, if constructed,
would have to encode increasingly fine information about the singular
conformal distribution.  It would not be an elementary asymptotic decay law
in the usual sense.

\section{Independent numerical checks and reproducibility}
\label{sec:numerics}

Numerics are not used to prove any exponent.  They are useful for detecting
normalization mistakes and for checking that the constants attached to the
proved formulas are consistent.

\subsection{Direct dyadic-cell quadrature}

Every zero of $P_n$ lies on the grid $2^{-n}\IntegerNumbers$.  On each of the $2^n$ finest
dyadic cells the absolute product is smooth, so fixed-order Gauss--Legendre
quadrature is highly effective.  Let
\begin{equation}
 I_n:=\int_0^1P_n(t)\Differential t,
 \qquad
 \mathcal A_n:=\rho^{-n}\int_0^1W_1(1+t)P_n(t)\Differential t.
 \label{eq:numeric-sequences}
\end{equation}
The computations below used eight Gauss points per finest cell and sixty tail
factors in $H$; changing to six, twelve, or sixteen points leaves the displayed
digits stable.

\begin{center}
\begin{tabular}{@{}rcc@{}}
\toprule
$n$ & $I_n/I_{n-1}$ & $\mathcal A_n$\\
\midrule
8  & 0.661322173532878 & 0.091266135409708\\
10 & 0.661322575197187 & 0.091266124753463\\
12 & 0.661322601094757 & 0.091266124150951\\
14 & 0.661322602051094 & 0.091266124131594\\
16 & 0.661322602062189 & 0.091266124131543\\
\bottomrule
\end{tabular}
\end{center}
This independently reproduces
\eqref{eq:rho-numeric} and \eqref{eq:A1}.  The analogous logarithmic-shell
sequence converges to
\begin{equation}
 B_1=0.065059235040389\ldots,
 \qquad B_1/a=0.093860635756799\ldots .
\end{equation}

A minimal implementation is:
\begin{lstlisting}[style=code,language=Python]
import numpy as np
from numpy.polynomial.legendre import leggauss

rho = 0.661322602060565
loga = np.log(2.0)
kappa = 0.5 + np.log(np.pi / rho) / loga


def H(y, terms=60):
    ans = np.ones_like(y)
    for m in range(1, terms + 1):
        ans *= np.sinc(y / 2**m)  # np.sinc(u)=sin(pi*u)/(pi*u)
    return np.abs(ans)


def shell_data(n, order=8):
    xi, wi = leggauss(order)
    total_I = 0.0
    total_A = 0.0
    cells = 2**n
    for j in range(cells):
        t = (j + (xi + 1.0)/2.0) / cells
        w = wi / (2.0*cells)
        P = np.ones_like(t)
        for ell in range(n):
            P *= np.abs(np.sin(np.pi * 2**ell * t))
        y = 1.0 + t
        W = (H(y) * (rho/np.pi) * y**(kappa-1.0)
             * np.exp(np.log(y)**2/(2.0*loga)))
        total_I += np.dot(w, P)
        total_A += np.dot(w, W*P)
    return total_I, total_A/rho**n
\end{lstlisting}

\subsection{The phase profile has narrow singular features}

At finite level define
\begin{equation}
 \mathscr F_n(z):=
 \frac{\int_0^zW_1(1+t)P_n(t)\Differential t}
      {\int_0^1W_1(1+t)P_n(t)\Differential t},
 \qquad
 \mathscr P_n(1+z):=\frac{1+\mathscr F_n(z)}{1+z}.
 \label{eq:finite-profile}
\end{equation}
For $n=13$, the same dyadic-cell quadrature gives
\begin{center}
\begin{tabular}{@{}ccc@{}}
\toprule
$y$ & $\mathscr P_{13}(y)$ & comment\\
\midrule
$8/7$  & 0.9163240341 & narrow lower feature\\
$1.150$ & 0.9248812037 & point seen on a grid of spacing $0.025$\\
$10/7$ & 1.1181043233 & narrow upper feature\\
$1.425$ & 1.1130801816 & point seen on a grid of spacing $0.025$\\
\bottomrule
\end{tabular}
\end{center}
At $n=17$, the values at $8/7$ and $10/7$ are approximately
$0.9162741$ and $1.1181341$.  These computations explain the discrepancy with
the comparative audit's coarser quoted range.  They are not certified global
minimum and maximum values: the limit is a singular-continuous distribution,
and ever finer features are expected.

\subsection{Finite-shell check of the Lambert formula}

Brute-force enumeration of the $2^n$ smooth lobe cells, followed by local
one-dimensional maximization, gives the following comparison.  The column
``Lambert core'' is the right side of \eqref{eq:annular-compact} without its
remainder.

\begin{center}
\begin{tabular}{@{}rccc@{}}
\toprule
$k$ & maximizing $t$ in $x=2^k(1+t)$ & $\log M_k$ &
$\log M_k-$Lambert core\\
\midrule
8  & 0.16637624 & $-39.27401113$  & $-2.00423$\\
12 & 0.16664866 & $-74.15774249$  & $-2.07119$\\
16 & 0.08333446 & $-119.95217544$ & $-2.04250$\\
\bottomrule
\end{tabular}
\end{center}
The bounded-looking residual at these modest levels is fully compatible with
the proved $\BigO((\log\log k)^2)$ error.  The jumps in the maximizing
mantissa reflect changes among preimages of the maximizing two-cycle.

\subsection{An exact-arithmetic certificate for a coarse Perron-root interval}

For completeness, the fourth report's Collatz--Wielandt certificate can be
checked with exact rational arithmetic.  Take
$h(x)=F(\sin\pi x)$, where
\[
 F(u)=1+\frac{376189}{10^6}u-\frac{69093}{10^6}u^2
      +\frac{15483}{10^6}u^3.
\]
The following Sturm-sequence check proves
$0.66126807<\rho<0.66134891$ without floating-point assumptions.

\begin{lstlisting}[style=code,language=Python]
import sympy as sp

t = sp.symbols("t", real=True)
u = 2*t/(1+t**2)
v = (1-t**2)/(1+t**2)

c = sp.Rational(376189, 10**6)
d = sp.Rational(-69093, 10**6)
e = sp.Rational(15483, 10**6)
F = lambda z: 1 + c*z + d*z**2 + e*z**3

R = sp.cancel((u*F(u) + v*F(v)) / (2*F(2*u*v)))
lo = sp.Rational(66126807, 10**8)
hi = sp.Rational(66134891, 10**8)

for expr in (R-lo, hi-R):
    num, den = map(lambda p: sp.Poly(p, t),
                   sp.together(expr).as_numer_denom())
    assert num.eval(0) > 0 and den.eval(0) > 0
    assert num.count_roots(0, 1) == 0
    assert den.count_roots(0, 1) == 0
\end{lstlisting}
The substitution $t=\tan(\pi x/4)$ converts
$(\LaplaceTransformSymbol_1h)(x)/h(x)$ into the rational function $R(t)$.  The two inequalities
then follow from Collatz--Wielandt.

\section{Conclusion}
\label{sec:conclusion}

The final answer is not one corrected equivalent.  It is a hierarchy generated
by the exact identity \eqref{eq:master}.

The deterministic denominator contributes
$-(\log x)^2/(2\log2)$.  The lacunary numerator then chooses the power of $x$:
$1/2$ for the geometric mean and almost-everywhere ray, $\rho_1$ for the
absolute mean, $1/\sqrt2$ for RMS, and $\sqrt3/2$ for the fixed-ray supremum.
The typical logarithm has genuine Gaussian and LIL fluctuations, with corrected
variance $\pi^2/4$.  The ordinary Ces\`aro question is governed by the
arithmetic exponent $\kappa_1$, not by the RMS exponent.  Its weaker bounded
form has the explicit analytic solution $A_1E_{\kappa_1}$; the all-$t$ limit
for that entire natural class is obstructed by a singular dyadic phase.  The
logarithmic mean removes that phase.

At the extremal end, the maximizing period-two orbit proves the sharp global
power exponent, while the largest value relative to a dyadic shell's left
endpoint is a different boundary-layer problem.  Its transition time is
selected by Lambert $W$ and creates a second log-normal correction in
$\log\log x$.  Finally, the local peak sequence retains the exact two-adic
multiplicity structure of the integer zeros, producing troughs far below the
annular maximum.

For a future Lean formalization, the development naturally separates into
independent blocks: locally uniform product convergence and dilation; the
canonical zero product and strict log-concavity of lobes; exact dyadic
factorization; the finite $L^2$ identity; the elementary two-step Gelfond
inequality; rational preimages and Lambert optimization; the boundary maximum
lemma; and, as imported analytic interfaces, Birkhoff/martingale limit theory
and the spectral theorem for $\LaplaceTransformSymbol_q$.  No ambiguity about ``the decay rate''
remains once the functional being normalized is stated.

\begin{thebibliography}{99}

\bibitem{ahl2017}
C.~Aistleitner, R.~Hofer, and G.~Larcher,
\newblock \emph{On evil Kronecker sequences and lacunary trigonometric products},
\newblock Ann. Inst. Fourier (Grenoble) \textbf{67} (2017), no.~2, 637--687;
preprint title \emph{On parametric Thue--Morse sequences and lacunary
trigonometric products}, arXiv:1502.06738,
\newblock \url{https://arxiv.org/abs/1502.06738}.

\bibitem{arias1982}
J.~Arias de Reyna,
\newblock \emph{An infinitely differentiable function with compact support:
Definition and properties},
\newblock Rev. Real Acad. Cienc. Exactas Fis. Natur. Madrid \textbf{76} (1982),
21--38; English translation, arXiv:1702.05442 (2017),
\newblock \url{https://arxiv.org/abs/1702.05442}.

\bibitem{baladi}
V.~Baladi,
\newblock \emph{Positive Transfer Operators and Decay of Correlations},
\newblock Advanced Series in Nonlinear Dynamics, vol.~16, World Scientific,
2000.

\bibitem{corless}
R.~M.~Corless, G.~H.~Gonnet, D.~E.~G.~Hare, D.~J.~Jeffrey, and D.~E.~Knuth,
\newblock \emph{On the Lambert $W$ function},
\newblock Adv. Comput. Math. \textbf{5} (1996), 329--359,
\newblock \url{https://doi.org/10.1007/BF02124750}.

\bibitem{fanschmeling}
A.~Fan, J.~Schmeling, and W.~Shen,
\newblock \emph{$L^\infty$-estimation of generalized Thue--Morse trigonometric
polynomials and ergodic maximization},
\newblock Discrete Contin. Dyn. Syst. \textbf{41} (2021), 297--327,
\newblock \url{https://doi.org/10.3934/dcds.2020363}.

\bibitem{fouvrymauduit}
\'E.~Fouvry and C.~Mauduit,
\newblock \emph{Sommes des chiffres et nombres presque premiers},
\newblock Math. Ann. \textbf{305} (1996), 571--599.

\bibitem{gelfond}
A.~O.~Gel'fond,
\newblock \emph{Sur les nombres qui ont des propri\'et\'es additives et
multiplicatives donn\'ees},
\newblock Acta Arith. \textbf{13} (1967/68), 259--265,
\newblock \url{https://doi.org/10.4064/aa-13-3-259-265}.

\bibitem{hallheyde}
P.~Hall and C.~C.~Heyde,
\newblock \emph{Martingale Limit Theory and Its Application},
\newblock Academic Press, New York, 1980.

\bibitem{mse-source}
V.~Reshetnikov,
\newblock \emph{Asymptotic decay rate of an infinite product of sinc functions},
\newblock Mathematics Stack Exchange, question 2745497 (2018),
\newblock \url{https://math.stackexchange.com/q/2745497}.

\bibitem{proveit}
V.~Reshetnikov,
\newblock \emph{ProveIt: Analysis/FabiusFunction}, GitHub repository,
\newblock \url{https://github.com/VladimirReshetnikov/ProveIt/tree/main/Analysis/FabiusFunction}.

\bibitem{ruelle}
D.~Ruelle,
\newblock \emph{Thermodynamic Formalism: The Mathematical Structures of
Classical Equilibrium Statistical Mechanics},
\newblock Addison--Wesley, 1978; second edition, Cambridge University Press,
2004.

\bibitem{stout}
W.~F.~Stout,
\newblock \emph{A martingale analogue of Kolmogorov's law of the iterated
logarithm},
\newblock Z. Wahrscheinlichkeitstheorie verw. Gebiete \textbf{15} (1970),
279--290.

\bibitem{wikipedia}
Wikipedia contributors,
\newblock \emph{Fabius function}, section ``Fourier transform'',
\newblock consulted 26 August 2026,
\newblock \url{https://en.wikipedia.org/wiki/Fabius_function}.

\bibitem{wave1a}
First-wave report 1,
\newblock \emph{Decay of the Fourier Transform of Rvachev's Up-Function:
Dyadic products, ergodic fluctuations, and sharp annular envelopes},
\newblock \texttt{ProveIt} research-frontier draft, 26 August 2026,
\newblock \url{https://github.com/VladimirReshetnikov/ProveIt/blob/main/Analysis/FabiusFunction/docs/semi-formalized-research-frontiers/drafts/rvachev_up_fourier_decay/rvachev_up_fourier_decay-1/rvachev_up_fourier_decay.tex}.

\bibitem{wave1b}
First-wave report 2,
\newblock \emph{Fourier Decay of Rvachev's Up-Function},
\newblock \texttt{ProveIt} research-frontier draft, 26 August 2026,
\newblock \url{https://github.com/VladimirReshetnikov/ProveIt/blob/main/Analysis/FabiusFunction/docs/semi-formalized-research-frontiers/drafts/rvachev_up_fourier_decay/Rvachev_Up_Fourier_Decay-2/Rvachev_Up_Fourier_Decay.tex}.

\bibitem{wave1c}
First-wave report 3,
\newblock \emph{Fourier Decay of Rvachev's Up-Function: Dyadic Dynamics,
Sharp Envelopes, and Mean Normalization},
\newblock \texttt{ProveIt} research-frontier draft, 26 August 2026,
\newblock \url{https://github.com/VladimirReshetnikov/ProveIt/blob/main/Analysis/FabiusFunction/docs/semi-formalized-research-frontiers/drafts/rvachev_up_fourier_decay/Rvachev_Up_Fourier_Decay-3/Fourier_decay_of_Rvachev_up.tex}.

\bibitem{wave1d}
First-wave report 4,
\newblock \emph{The Fourier Decay of Rvachev's Up-Function},
\newblock \texttt{ProveIt} research-frontier draft, 26 August 2026,
\newblock \url{https://github.com/VladimirReshetnikov/ProveIt/blob/main/Analysis/FabiusFunction/docs/semi-formalized-research-frontiers/drafts/rvachev_up_fourier_decay/rvachev_up_fourier_decay-4/rvachev_up_fourier_decay.tex}.

\bibitem{comparative}
\emph{Fourier Decay of Rvachev's Up-Function: Comparative Audit and Resolution},
\newblock \texttt{ProveIt} research-frontier draft, 26 August 2026,
\newblock \url{https://github.com/VladimirReshetnikov/ProveIt/blob/main/Analysis/FabiusFunction/docs/semi-formalized-research-frontiers/drafts/rvachev_up_fourier_decay/Rvachev_Up_Fourier_Decay_Comparative_Audit/Rvachev_Up_Fourier_Decay_Comparative_Audit.tex}.

\end{thebibliography}

\end{document}
